\chapter{La lecture}

\section{Aléatoire et répétition}
La lecture aléatoire choisit le morceau suivant au hasard dans la liste ou le dossier. La répétition rejoue le
morceau en cours sans fin ; sans elle, le morceau s'arrête à la fin de sa durée connue.

\section{Durée, boucles et fondus}
Beaucoup de musiques de jeux bouclent indéfiniment : le fichier ne dit pas quand s'arrêter. Quand une durée est
connue, par le format lui-même ou par une base de données (la base \emph{Songlengths} de la HVSC pour le SID), nous
l'utilisons. Sinon, nous appliquons la durée par défaut, que vous réglez dans les réglages, ainsi que la limite de
boucles et les fondus. Quand un morceau a une introduction avant sa boucle, la barre de progression marque le début
de la partie bouclée.

\section{Hors de l'application}
Vous pouvez piloter la lecture depuis la notification, l'écran verrouillé, un casque Bluetooth ou Android Auto.
Android Auto affiche vos listes et une visualisation simplifiée.

\section{Partager}
Le menu du morceau partage le fichier, envoie un lien vers le morceau, crée une sonnerie à partir d'un extrait et
affiche toutes les propriétés connues du morceau.
